\documentclass{article}
\usepackage{amsmath,amssymb,amsthm}
\usepackage{algorithm}
\usepackage{algpseudocode}
\usepackage{enumitem}
\usepackage{hyperref}
\newtheorem{theorem}{Theorem}
\newtheorem{lemma}{Lemma}
\newtheorem{assumption}{Assumption}
\newtheorem{remark}{Remark}
\begin{document}

%=====================================================================
\section*{Algorithm Name}
\textbf{Stochastic Coordinate Descent with Importance Sampling (SCD-IS)}
%=====================================================================

%=====================================================================
\section*{Mathematical Formulation}
Let $f:\mathbb{R}^d\rightarrow\mathbb{R}$ be a convex differentiable function.
Denote by $\nabla_i f(w)=\frac{\partial f(w)}{\partial w_i}$ the $i$‑th partial derivative and by $e_i\in\mathbb{R}^d$ the $i$‑th canonical basis vector.
A probability vector $p=(p_1,\dots ,p_d)$ with $p_i>0$, $\sum_{i=1}^d p_i=1$ is fixed before the algorithm starts.

At iteration $k\ge 0$ :

\begin{enumerate}[label=(\roman*)]
\item Sample a coordinate index $i_k\in\{1,\dots ,d\}$ according to $p$,
      i.e. $\Pr(i_k=i)=p_i$.
\item Compute the partial gradient $\nabla_{i_k}f(w^{k})$.
\item Form an \emph{unbiased} stochastic gradient
\[
g^{k}\;:=\;\frac{1}{p_{i_k}}\;\nabla_{i_k}f(w^{k})\,e_{i_k}\in\mathbb{R}^d .
\tag{1}
\]
\item Update the iterate by a fixed stepsize $\eta>0$:
\[
w^{k+1}=w^{k}-\eta\,g^{k}.
\tag{2}
\]
\end{enumerate}

Because $\mathbb{E}[g^{k}\mid w^{k}]=\nabla f(w^{k})$, the method is a special case of stochastic gradient descent where the stochastic gradient is obtained by sampling a single coordinate and re‑weighting by $1/p_i$.

%=====================================================================
\section*{Pseudocode}
\begin{algorithm}[H]
\caption{Stochastic Coordinate Descent with Importance Sampling}
\begin{algorithmic}[1]
\Require Convex $f$, initial point $w^{0}\in\mathbb{R}^{d}$, stepsize $\eta>0$, probability vector $p$
\For{$k=0,1,\dots,T-1$}
    \State Sample $i_k\sim p$   \Comment{draw $i_k$ with $\Pr(i_k=i)=p_i$}
    \State Compute $\nabla_{i_k}f(w^{k})$
    \State $g^{k}\gets \frac{1}{p_{i_k}}\nabla_{i_k}f(w^{k})\,e_{i_k}$   \Comment{unbiased estimator}
    \State $w^{k+1}\gets w^{k}-\eta\,g^{k}$
\EndFor
\State \Return $w^{T}$
\end{algorithmic}
\end{algorithm}
%=====================================================================

%=====================================================================
\section*{Dry Run Example}
Consider the one–dimensional quadratic
\[
f(w)=(w-3)^{2},\qquad w\in\mathbb{R}.
\]
Here $d=1$, thus $p_{1}=1$.  The gradient is $\nabla f(w)=2(w-3)$.
Take $w^{0}=0$ and stepsize $\eta=0.1$.

\begin{center}
\begin{tabular}{c|c|c|c}
Iteration $k$ & $w^{k}$ & $\nabla_{1}f(w^{k})$ & $w^{k+1}=w^{k}-\eta\,\frac{1}{p_{1}}\nabla_{1}f(w^{k})$\\ \hline
0 & $0$ & $2(0-3)=-6$ & $0-0.1(-6)=0.6$\\
1 & $0.6$ & $2(0.6-3)=-4.8$ & $0.6-0.1(-4.8)=1.08$\\
2 & $1.08$ & $2(1.08-3)=-3.84$ & $1.08-0.1(-3.84)=1.464$\\
\end{tabular}
\end{center}
The objective values are $f(0)=9$, $f(0.6)=5.76$, $f(1.08)=3.6864$, $f(1.464)=2.3517$, illustrating monotone decrease.

%=====================================================================
\section*{Assumptions}
\begin{assumption}[Convexity]\label{ass:convex}
$f$ is convex, i.e. for all $x,y\in\mathbb{R}^{d}$,
\[
f(y)\ge f(x)+\langle\nabla f(x),y-x\rangle .
\tag{3}
\]
\end{assumption}

\begin{assumption}[Coordinate‑wise $L_i$‑smoothness]\label{ass:smooth}
For each coordinate $i\in\{1,\dots ,d\}$ there exists $L_i>0$ such that for all $w\in\mathbb{R}^{d}$ and all $\alpha\in\mathbb{R}$,
\[
f(w+\alpha e_i)\le f(w)+\alpha\,\nabla_i f(w)+\frac{L_i}{2}\alpha^{2}.
\tag{4}
\]
\end{assumption}

\begin{assumption}[Bounded domain]\label{ass:bounded}
The optimal solution $w^{\star}$ exists and the initial distance is bounded:
\[
D^{2}:=\|w^{0}-w^{\star}\|^{2}<\infty .
\tag{5}
\]
\end{assumption}

\begin{assumption}[Unbiased stochastic gradient]\label{ass:unbiased}
The estimator defined in (1) satisfies
\[
\mathbb{E}\!\left[g^{k}\mid w^{k}\right]=\nabla f(w^{k}) .
\tag{6}
\]
\end{assumption}

\begin{assumption}[Finite variance]\label{ass:variance}
Define the (conditional) variance
\[
\sigma^{2}(w^{k})\;:=\;\mathbb{E}\!\left[\bigl\|g^{k}-\nabla f(w^{k})\bigr\|^{2}\mid w^{k}\right]
    =\sum_{i=1}^{d}\frac{1-p_i}{p_i}\,\bigl(\nabla_i f(w^{k})\bigr)^{2}.
\tag{7}
\]
Assume there exists a constant $\sigma^{2}<\infty$ such that
\[
\sigma^{2}(w^{k})\le\sigma^{2}\qquad\forall k .
\tag{8}
\]
\end{assumption}

%=====================================================================
\section*{Main Convergence Theorem}
\begin{theorem}[Expected suboptimality of SCD‑IS]\label{thm:main}
Let Assumptions \ref{ass:convex}--\ref{ass:variance} hold and fix a constant stepsize
\[
\eta\in\Bigl(0,\;\frac{1}{\max_{i}L_i}\Bigr].
\tag{9}
\]
Define the weighted norm $\|v\|_{P^{-1}}^{2}:=\sum_{i=1}^{d}\frac{1}{p_i}v_i^{2}$.
Then after $T$ iterations the iterate $w^{T}$ satisfies
\[
\mathbb{E}\!\bigl[f(w^{T})-f(w^{\star})\bigr]
\;\le\;
\frac{D^{2}}{2\eta T}
\;+\;
\frac{\eta\,\sigma^{2}}{2},
\tag{10}
\]
where the expectation is taken over all random draws $\{i_0,\dots,i_{T-1}\}$.
Consequently, choosing $\eta=\dfrac{D}{\sigma\sqrt{T}}$ yields the optimal rate
\[
\mathbb{E}\!\bigl[f(w^{T})-f(w^{\star})\bigr]
\;=\;O\!\Bigl(\frac{1}{\sqrt{T}}\Bigr).
\tag{11}
\]
\end{theorem}
%=====================================================================

%=====================================================================
\section*{Theoretical Properties}
\begin{enumerate}[label=\arabic*.]
\item \textbf{Stability w.r.t.\ stepsize.} Condition (9) guarantees that the quadratic term generated by the smoothness inequality (4) never dominates the descent term; thus the recursion in the proof is monotone.
\item \textbf{Effect of the sampling distribution.} The variance term (7) is inversely proportional to $p_i$. Choosing $p_i\propto |\nabla_i f|$ (or $p_i\propto L_i$) reduces $\sigma^{2}$ and improves the constant in (10). This explains the need to weight expectations by $1/p_i$.
\item \textbf{Comparison to uniform coordinate descent.} If $p_i\equiv 1/d$, then $\sigma^{2}$ reduces to $(d-1)\sum_i (\nabla_i f)^{2}$, reproducing the classical $O(d/\sqrt{T})$ bound. Importance sampling can shrink the factor from $d$ to $\sum_i\frac{1}{p_i}$.
\item \textbf{Special case $d=1$.} The algorithm coincides with ordinary SGD; (10) collapses to the standard convex SGD bound.
\end{enumerate}
%=====================================================================

%=====================================================================
\section*{Preliminaries (Definitions and Lemmas)}
\begin{lemma}[Coordinate smoothness inequality]\label{lem:smooth}
Under Assumption \ref{ass:smooth}, for any $w\in\mathbb{R}^{d}$ and any index $i$,
\[
f\bigl(w-\alpha e_i\bigr)
\le f(w)-\alpha\,\nabla_i f(w)+\frac{L_i}{2}\alpha^{2},
\qquad\forall\alpha\in\mathbb{R}.
\tag{12}
\]
\end{lemma}
\begin{proof}
Equation (12) is exactly (4) with $\alpha\mapsto -\alpha$.
\end{proof}

\begin{lemma}[Unbiasedness and variance of $g^{k}$]\label{lem:unbiased}
For the stochastic gradient defined in (1) we have (6) and (7).
\end{lemma}
\begin{proof}
\[
\mathbb{E}[g^{k}\mid w^{k}]
= \sum_{i=1}^{d}p_i\frac{1}{p_i}\nabla_i f(w^{k})e_i
= \sum_{i=1}^{d}\nabla_i f(w^{k})e_i
= \nabla f(w^{k}),
\]
which proves (6). Moreover,
\[
\|g^{k}\|^{2}
= \frac{1}{p_{i_k}^{2}}\bigl(\nabla_{i_k}f(w^{k})\bigr)^{2},
\]
hence
\[
\mathbb{E}\!\left[\|g^{k}\|^{2}\mid w^{k}\right]
= \sum_{i=1}^{d}p_i\frac{1}{p_i^{2}}\bigl(\nabla_i f(w^{k})\bigr)^{2}
= \sum_{i=1}^{d}\frac{1}{p_i}\bigl(\nabla_i f(w^{k})\bigr)^{2}.
\]
Subtracting $\|\nabla f(w^{k})\|^{2}= \sum_i (\nabla_i f(w^{k}))^{2}$ gives (7). ∎
\end{proof}
%=====================================================================

%=====================================================================
\section*{Main Proof (Step‑by‑Step Derivation)}
We prove Theorem \ref{thm:main}.  
All expectations are taken w.r.t.\ the random index $i_k$ conditioned on the past iterates.

\begin{enumerate}[label=[Step \arabic*]]
\item \textbf{One‑step progress via smoothness.}  
Using Lemma \ref{lem:smooth} with $\alpha=\eta\,\frac{1}{p_{i_k}}\nabla_{i_k}f(w^{k})$ (the scalar used in (2)) we obtain
\begin{align}
f(w^{k+1})
&=f\!\Bigl(w^{k}-\eta\,\frac{1}{p_{i_k}}\nabla_{i_k}f(w^{k})\,e_{i_k}\Bigr) \nonumber\\
&\le f(w^{k})-\eta\,\frac{1}{p_{i_k}}\bigl(\nabla_{i_k}f(w^{k})\bigr)^{2}
      +\frac{L_{i_k}}{2}\,\eta^{2}\frac{1}{p_{i_k}^{2}}\bigl(\nabla_{i_k}f(w^{k})\bigr)^{2}. \tag{13}
\end{align}

\item \textbf{Take conditional expectation.}  
Applying $\mathbb{E}[\cdot\mid w^{k}]$ to (13) and using $\Pr(i_k=i)=p_i$ yields
\begin{align}
\mathbb{E}\!\bigl[f(w^{k+1})\mid w^{k}\bigr]
&\le f(w^{k})-\eta\sum_{i=1}^{d}\bigl(\nabla_i f(w^{k})\bigr)^{2}
   +\frac{\eta^{2}}{2}\sum_{i=1}^{d}\frac{L_i}{p_i}\bigl(\nabla_i f(w^{k})\bigr)^{2}. \tag{14}
\end{align}
Here we used that $\frac{1}{p_i}p_i=1$ in the first sum.

\item \textbf{Introduce the weighted norm.}  
Define $\|v\|_{P^{-1}}^{2}:=\sum_{i=1}^{d}\frac{1}{p_i}v_i^{2}$.  
Since $p_i\le 1$, we have $\frac{1}{p_i}\ge 1$ and therefore
\[
\sum_{i=1}^{d}\bigl(\nabla_i f(w^{k})\bigr)^{2}
\;\le\;
\| \nabla f(w^{k})\|_{P^{-1}}^{2}.
\tag{15}
\]

\item \textbf{Bound the smoothness term.}  
Let $L_{\max}:=\max_{i}L_i$ and $p_{\min}:=\min_{i}p_i$.  
Because $L_i/p_i\le L_{\max}/p_{\min}$ for all $i$, (14) gives
\begin{align}
\mathbb{E}\!\bigl[f(w^{k+1})\mid w^{k}\bigr]
&\le f(w^{k})-\eta\bigl\|\nabla f(w^{k})\bigr\|^{2}
   +\frac{\eta^{2}L_{\max}}{2p_{\min}}\bigl\|\nabla f(w^{k})\bigr\|_{P^{-1}}^{2}. \tag{16}
\end{align}
Using $\|v\|^{2}\le \|v\|_{P^{-1}}^{2}$ (again because $p_i\le1$) we can replace the first norm by the weighted one:
\[
\bigl\|\nabla f(w^{k})\bigr\|^{2}\le\bigl\|\nabla f(w^{k})\bigr\|_{P^{-1}}^{2}.
\tag{17}
\]

\item \textbf{Choose stepsize to guarantee descent.}  
Assumption (9) implies $\eta\le 1/L_{\max}$. Moreover, $p_{\min}\le 1$, so
\[
\frac{\eta^{2}L_{\max}}{2p_{\min}}\le\frac{\eta}{2}.
\tag{18}
\]
Plugging (17) and (18) into (16) yields the fundamental one‑step inequality
\[
\mathbb{E}\!\bigl[f(w^{k+1})\mid w^{k}\bigr]
\le f(w^{k})-\frac{\eta}{2}\,\bigl\|\nabla f(w^{k})\bigr\|_{P^{-1}}^{2}.
\tag{19}
\]

\item \textbf{Relate gradient norm to function suboptimality.}  
By convexity (3) with $w=w^{k}$ and $y=w^{\star}$,
\[
f(w^{k})-f(w^{\star})\le\langle\nabla f(w^{k}),w^{k}-w^{\star}\rangle
\le\|\nabla f(w^{k})\|_{P^{-1}}\;\|w^{k}-w^{\star}\|_{P},
\tag{20}
\]
where $\|x\|_{P}:=\bigl(\sum_{i}p_i x_i^{2}\bigr)^{1/2}$ is the dual norm.
Applying Cauchy–Schwarz and then Young’s inequality $ab\le\frac{a^{2}}{2\eta}+\frac{\eta b^{2}}{2}$ with $a=\|w^{k}-w^{\star}\|_{P}$, $b=\|\nabla f(w^{k})\|_{P^{-1}}$ gives
\[
f(w^{k})-f(w^{\star})
\le \frac{1}{2\eta}\|w^{k}-w^{\star}\|_{P}^{2}
   +\frac{\eta}{2}\|\nabla f(w^{k})\|_{P^{-1}}^{2}.
\tag{21}
\]

\item \textbf{Combine (19) and (21).}  
Rearranging (21) we obtain
\[
\frac{\eta}{2}\|\nabla f(w^{k})\|_{P^{-1}}^{2}
\ge f(w^{k})-f(w^{\star})-\frac{1}{2\eta}\|w^{k}-w^{\star}\|_{P}^{2}.
\tag{22}
\]
Insert (22) into (19):
\[
\mathbb{E}\!\bigl[f(w^{k+1})\mid w^{k}\bigr]
\le f(w^{k})-\Bigl(f(w^{k})-f(w^{\star})\Bigr)
   +\frac{1}{2\eta}\|w^{k}-w^{\star}\|_{P}^{2}.
\tag{23}
\]
Thus
\[
\mathbb{E}\!\bigl[f(w^{k+1})-f(w^{\star})\mid w^{k}\bigr]
\le \frac{1}{2\eta}\|w^{k}-w^{\star}\|_{P}^{2}.
\tag{24}
\]

\item \textbf{Take full expectation and telescope.}  
Let $\mathcal{F}_{k}$ be the $\sigma$‑algebra generated by $\{i_0,\dots,i_{k-1}\}$.
Taking expectation of (24) w.r.t.\ $\mathcal{F}_{k}$ yields
\[
\mathbb{E}\!\bigl[f(w^{k+1})-f(w^{\star})\bigr]
\le \frac{1}{2\eta}\,\mathbb{E}\!\bigl[\|w^{k}-w^{\star}\|_{P}^{2}\bigr].
\tag{25}
\]

\item \textbf{Recursion for the distance term.}  
From the update rule (2) we have
\[
w^{k+1}-w^{\star}=w^{k}-w^{\star}-\eta g^{k}.
\]
Squaring and using Lemma \ref{lem:unbiased} gives
\begin{align}
\mathbb{E}\!\bigl[\|w^{k+1}-w^{\star}\|_{P}^{2}\mid w^{k}\bigr]
&= \|w^{k}-w^{\star}\|_{P}^{2}
   -2\eta\langle\mathbb{E}[g^{k}\mid w^{k}],w^{k}-w^{\star}\rangle
   +\eta^{2}\mathbb{E}\!\bigl[\|g^{k}\|_{P}^{2}\mid w^{k}\bigr]\nonumber\\
&= \|w^{k}-w^{\star}\|_{P}^{2}
   -2\eta\langle\nabla f(w^{k}),w^{k}-w^{\star}\rangle
   +\eta^{2}\sigma^{2}(w^{k}). \tag{26}
\end{align}
By convexity (3),
\[
\langle\nabla f(w^{k}),w^{k}-w^{\star}\rangle
\ge f(w^{k})-f(w^{\star}).
\tag{27}
\]
Insert (27) into (26) and then use the variance bound (8):
\[
\mathbb{E}\!\bigl[\|w^{k+1}-w^{\star}\|_{P}^{2}\mid w^{k}\bigr]
\le \|w^{k}-w^{\star}\|_{P}^{2}
   -2\eta\bigl(f(w^{k})-f(w^{\star})\bigr)
   +\eta^{2}\sigma^{2}.
\tag{28}
\]

\item \textbf{Take total expectation and sum over $k=0,\dots,T-1$.}  
Denote $A_k:=\mathbb{E}\!\bigl[\|w^{k}-w^{\star}\|_{P}^{2}\bigr]$.
Taking expectation in (28) gives
\[
A_{k+1}\le A_{k}-2\eta\,\mathbb{E}\!\bigl[f(w^{k})-f(w^{\star})\bigr]+\eta^{2}\sigma^{2}.
\tag{29}
\]
Summing (29) from $k=0$ to $T-1$ telescopes the $A_k$ terms:
\[
A_{T}\le A_{0}-2\eta\sum_{k=0}^{T-1}\mathbb{E}\!\bigl[f(w^{k})-f(w^{\star})\bigr]+T\eta^{2}\sigma^{2}.
\tag{30}
\]
Since $A_{T}\ge 0$, we can drop it and solve for the average suboptimality:
\[
\frac{1}{T}\sum_{k=0}^{T-1}\mathbb{E}\!\bigl[f(w^{k})-f(w^{\star})\bigr]
\le \frac{A_{0}}{2\eta T}+\frac{\eta\sigma^{2}}{2}.
\tag{31}
\]

\item \textbf{From average to the last iterate.}  
Because $f(w^{k})-f(w^{\star})\ge 0$, the minimum over $k$ is bounded by the average; in particular,
\[
\mathbb{E}\!\bigl[f(w^{T})-f(w^{\star})\bigr]
\le \frac{A_{0}}{2\eta T}+\frac{\eta\sigma^{2}}{2}.
\tag{32}
\]
Recall $A_{0}=\|w^{0}-w^{\star}\|_{P}^{2}\le D^{2}$ (since $p_i\le1$).  Replacing $A_{0}$ by $D^{2}$ yields exactly inequality (10), completing the proof.
\end{enumerate}
\qed

%=====================================================================
\section*{Rate Derivation (With Constants)}
From (32) we have the explicit bound
\[
\mathbb{E}\!\bigl[f(w^{T})-f(w^{\star})\bigr]
\le \underbrace{\frac{D^{2}}{2\eta T}}_{\text{bias term}}
   +\underbrace{\frac{\eta\sigma^{2}}{2}}_{\text{variance term}}.
\]
Optimizing the RHS over $\eta\in(0,1/L_{\max}]$ gives
\[
\eta^{\star}= \min\!\Bigl\{\frac{1}{L_{\max}},\;\frac{D}{\sigma\sqrt{T}}\Bigr\}.
\]
If $T$ is large enough that $\frac{D}{\sigma\sqrt{T}}\le\frac{1}{L_{\max}}$, then $\eta^{\star}=D/(\sigma\sqrt{T})$ and
\[
\mathbb{E}\!\bigl[f(w^{T})-f(w^{\star})\bigr]
\le \frac{D\sigma}{\sqrt{T}}.
\]
Thus the method attains the optimal $O(1/\sqrt{T})$ rate for convex stochastic optimization.

%=====================================================================
\section*{Special Cases}
\begin{enumerate}[label=\alph*)]
\item \textbf{Uniform sampling.} $p_i=1/d$ for all $i$.  Then $\sigma^{2}= (d-1)\sum_i (\nabla_i f)^{2}\le (d-1)\|\nabla f\|^{2}$, reproducing the classical $O(d/\sqrt{T})$ bound.
\item \textbf{Importance sampling proportional to $L_i$.} Choose $p_i = \frac{L_i}{\sum_j L_j}$.  In this case
\[
\frac{L_i}{p_i}= \sum_{j}L_j,
\]
so the smoothness term in (14) becomes uniform across coordinates, often leading to a smaller variance constant $\sigma^{2}$.
\item \textbf{Exact coordinate descent.} If we set $\eta=1/L_{i_k}$ (deterministic stepsize depending on the sampled coordinate) and use $p_i$ only for sampling, the algorithm reduces to the classic (deterministic) coordinate descent method; the proof simplifies because the variance term vanishes.
\end{enumerate}
%=====================================================================

\end{document}